% !TeX program = xelatex
\documentclass[11pt,a4paper]{article}
\usepackage[margin=2.4cm]{geometry}
\usepackage{fontspec}
\usepackage{polyglossia}
\setdefaultlanguage{vietnamese}
\setmainfont{Arial}[Path=../assets/fonts/,Extension=.ttf,
    UprightFont=*-Regular,BoldFont=*-Bold,ItalicFont=*-Italic,BoldItalicFont=*-BoldItalic]
\usepackage{amsmath}
\usepackage[hidelinks]{hyperref}
\setlength{\parindent}{0pt}
\setlength{\parskip}{0.55em}
\title{Ghi chú thuyết trình Bootstrap}
\author{Nhóm D}
\date{}
\begin{document}
\maketitle

\section{Cách trình bày trong 30 phút}
Dành 4 phút cho trang 1--4, 8 phút cho trang 5--9, 5 phút cho hai ví dụ
ở trang 10--11, 6 phút cho mô phỏng ở trang 12--15, 5 phút cho ứng dụng
và đổi thang ở trang 16--18, rồi 2 phút cho kết luận và cảm ơn.
Trang 4 đặt cạnh nhau mức đặt trước và kết quả đo trong mô phỏng.
Khi đến trang 13, chỉ vào chỉ báo và giải thích vì sao mỗi khoảng
chứa giá trị thật góp đúng một đơn vị vào tử số.

\section{Ký hiệu và các công thức cần giải thích}
\textbf{Ký hiệu.}
$n$ là số quan sát trong mẫu gốc và trong mỗi mẫu lấy lại.
$N$ là số bản sao Bootstrap trên một mẫu gốc. $j$ đánh số bản sao từ $1$ đến $N$,
$i$ đánh số vị trí từ $1$ đến $n$, còn $k$ đánh số bộ dữ liệu mô phỏng từ $1$ đến $K$.
Báo cáo và mã R viết $B$ cho $N$, và $b$ cho $j$. Đây là hai cách đặt tên cho
cùng một số lượng và chỉ số. Chữ $J$ trong bản cũ là chỉ số chọn vị trí
trong mẫu gốc. Có thể diễn đạt thao tác đó bằng lời và giữ $j$ cho số thứ tự bản sao.
Ký hiệu $\mathcal N(\mu,\sigma^2)$ chỉ phân phối chuẩn.

\textbf{Mức danh nghĩa và độ bao phủ.}
Gọi $C(X)=[L(X),U(X)]$ là khoảng dựng từ dữ liệu $X$.
Mức danh nghĩa là $1-\alpha$. Độ bao phủ của thủ tục là
$\Pr_\theta\{\theta\in C(X)\}$ khi lấy mẫu lặp lại ở cùng thiết lập.
Bài giảng Chương 4, trang 55, phát biểu điều kiện
$\Pr_\theta\{L(X)\le\theta\le U(X)\}\ge1-\alpha$ cho một khoảng ở mức đó.
Báo cáo nhóm D (Lời mở đầu, tr. 3 và nội dung chính, tr. 12--14) dùng cùng cách hiểu và đo độ bao phủ
của các khoảng Bootstrap xấp xỉ bằng mô phỏng. Con số $85{,}47\%$
là ước lượng thực nghiệm từ $K=10\,000$ lần, có sai số Monte Carlo.
Nếu muốn tính sai số chuẩn của tỷ lệ, đặt $Y_k=I(\theta\in C_k)$.
Vì $Y_k$ là biến 0--1, phương sai của nó là $CP(1-CP)$, nên
\[
    \widehat{MCSE}=\sqrt{\widehat{CP}(1-\widehat{CP})/K}
        \approx 0{,}00352
\]
cho Basic ở kịch bản mẫu mũ $n=12$, tương đương khoảng $0{,}35$ điểm phần trăm.

\textbf{Sai số chuẩn.}
Trong mô phỏng lấy lại ngẫu nhiên, dùng độ lệch chuẩn mẫu của $N$ bản sao
với mẫu số $N-1$. Trong hai ví dụ liệt kê đủ $n^n$ mẫu có thứ tự, ta biết
toàn bộ phân phối có điều kiện. Lúc đó tính phương sai theo xác suất chính
xác bằng tổng các bình phương độ lệch nhân với xác suất tương ứng.
Mỗi bộ chỉ số có cùng xác suất $1/n^n$, kể cả khi các vị trí đó mang giá trị trùng.
Với type 1, ở mức $p\in(0,1)$, lấy giá trị ở hạng $\lceil mp\rceil$
trong $m$ giá trị đã sắp xếp, có tính cả các giá trị lặp.

\textbf{Trung bình mẫu chuẩn.}
Nếu $X_1,\ldots,X_n$ độc lập, mỗi quan sát có phân phối $\mathcal N(\mu,\sigma^2)$,
thì trung bình mẫu cũng có phân phối chuẩn với kỳ vọng $\mu$ và phương sai
$\sigma^2/n$. Kỳ vọng theo tính tuyến tính, còn tính độc lập cho
$\operatorname{Var}(\bar X)=n\sigma^2/n^2$. Tính chuẩn đến từ tính đóng
của phân phối chuẩn đối với tổ hợp tuyến tính của các biến chuẩn độc lập.
Ý này giúp hiểu tổng thể chuẩn trong mô phỏng. Với mẫu nhỏ, thay sai số chuẩn
bằng ước lượng Bootstrap vẫn có thể cho độ bao phủ dưới mức danh nghĩa.

\textbf{Đổi thang log.}
Ở ví dụ trang 18, đặt $\eta=\log\theta$. Với mỗi bản sao $\hat\theta_j^*>0$,
tạo bản sao tương ứng $\hat\eta_j^*=\log(\hat\theta_j^*)$.
Nếu dùng type 1 thì hai phân vị trên thang log là log của hai phân vị gốc.
Với Basic, hai cách lấy cận dưới lần lượt cho
$\log(2\hat\theta-q_U^*)$ và $2\log\hat\theta-\log q_U^*$.
Chúng có thể khác nhau vì log không phân phối qua phép trừ.
Với phân vị nội suy type 7, log của một giá trị nội suy cũng có thể khác
giá trị nội suy của hai log.

\section{R1--R6 trong bản slide trước}
\begin{enumerate}
    \item R1 (ECDF). Sinh mẫu chuẩn ở nhiều cỡ mẫu, tính ECDF và so sánh
    với hàm phân phối chuẩn. Đầu ra là đồ thị đặt các ECDF cạnh đường phân phối lý thuyết.
    \item R2 (mẫu đối xứng). Liệt kê ba vị trí rút từ mẫu $(1,2,3)$, tính
    trung bình mỗi mẫu, đếm số lần xuất hiện của từng trung bình và lấy phân vị
    type 1. Các tần số trên slide 10 dùng cách tính này.
    \item R3 (dựng khoảng). Từ một bộ dữ liệu, tạo $N$ mẫu lấy lại.
    Tính độ lệch chuẩn của các bản sao, hai phân vị type 7 rồi áp dụng
    công thức Normal, Basic và Percentile. Đoạn này là phần tính cận
    trên một bộ dữ liệu trong vòng ngoài của mô phỏng.
    \item R4 (vẽ BCa). Nhập RV và các cận đã công bố, sau đó vẽ điểm
    ước lượng và đoạn biểu diễn khoảng. Việc tái vẽ sử dụng số trong bảng
    của bài báo. Dựng lại khoảng BCa từ đầu cần dữ liệu gốc.
    \item R5 (HR phân tầng). Sinh dữ liệu minh họa, lấy lại bệnh nhân
    trong từng nhóm và khớp lại mô hình Cox trên mỗi mẫu Bootstrap.
    Dãy HR thu được dùng để tính sai số chuẩn và phân vị. Các con số
    chọn sẵn trên slide 16 thuộc ví dụ giả định trong báo cáo.
    \item R6 (đổi thang log). Dùng cùng dãy trung bình Bootstrap cho hai
    thang. So sánh log của cận Basic gốc với cận Basic dựng trực tiếp
    từ các log. Trong mã R đầy đủ, phép đối chiếu còn chạy cho Normal
    và Percentile.
\end{enumerate}

\section{Đối chiếu tài liệu và nội dung đã rút gọn}
Chân trang của slide dùng số trang in trong \texttt{main.pdf} của nhóm,
với mục lục đánh số i--ii và nội dung chính bắt đầu lại từ trang 1 (trang thứ 7 của PDF). Với bài giảng Chương 4, dùng số
trang ghi trên bài giảng. Bài báo Deng và cộng sự được dẫn theo Bảng 1
trong \href{https://link.springer.com/article/10.1186/1477-7525-11-89}{bản toàn văn công khai}.
Các nguồn phụ và ghi chú nội bộ không cần cho phần trình chiếu đã được bỏ.

Bản 20 trang bỏ slide quy ước ký hiệu cũ và đưa các ký hiệu vào đúng chỗ dùng.
Các trang 34 và 36 cũ cùng phần phụ lục bảo vệ và các trang mã R được rút khỏi
bộ chiếu. Nội dung so sánh và công thức vẫn có ở phần chính hoặc ở ghi chú này.
Các ví dụ trên bộ chiếu có dữ kiện, câu hỏi và phép thế số. Phần ứng dụng
phân biệt HR giả định với RV lấy từ nghiên cứu đã công bố.

\end{document}
